% ─────────────────────────────────────────────────────────────────────────────
\section{Le registre incr\'emental}
% ─────────────────────────────────────────────────────────────────────────────

\begin{maitre}
Si le stock change l\'eg\`erement (ajout de deux produits),
faut-il recalculer toutes les port\'ees~?
\end{maitre}

\begin{apprenti}
Non, seulement pour les nouveaux produits.
Les produits existants n'ont pas chang\'e.
\end{apprenti}

\begin{maitre}
C'est exactement ce que fait le \textbf{registre de calcul}.
Il m\'emorise les paires (mat\'eriau, config) d\'ej\`a calcul\'ees
dans un fichier CSV persistant.
\end{maitre}

\begin{lstlisting}[style=python, caption={Registre incrémental --- extrait de \fichier{registre.py}}]
class RegistreCalcul:
    """Mémorise les calculs déjà effectués pour éviter les doublons."""

    def est_calcule(self, id_mat: str, id_calc: str) -> bool:
        """Retourne True si ce couple a déjà été calculé."""
        if self._force_recalcul:
            return False   # ignorer le registre si recalcul forcé
        return (id_mat, id_calc) in self._calcules

    def enregistrer(self, id_mat: str, id_calc: str, statut: str):
        """Marque ce couple comme calculé."""
        self._calcules.add((id_mat, id_calc))
\end{lstlisting}

Et dans la boucle principale du moteur~:

\begin{lstlisting}[style=python, caption={Usage du registre dans \fichier{moteur.py}}]
for config in app_config.configs_calcul:
    id_calc = config.id_config_calcul

    # Si déjà calculé -> sauter (sauf si --recalcul-complet)
    if not registre._force_recalcul and registre.est_calcule(id_mat, id_calc):
        if verbose:
            logger.info(f"Déjà calculé ({id_mat}, {id_calc}) - ignoré.")
        continue   # "continue" = passer à l'itération suivante

    # ... calculs EC5 ...

    registre.enregistrer(id_mat, id_calc, "calcule")
\end{lstlisting}

Pour forcer un recalcul complet~:

\begin{lstlisting}[style=toml, language={}, caption={Forcer le recalcul complet}]
uv run abac calculer --config config.toml --recalcul-complet
\end{lstlisting}
